\section{Technical Appendix and Deep Dives}

\begin{frame}[label=backup_rootsystem]{Formal Definition: Root System}
    \textbf{Definition 4.4.1:} Let $E$ be a finite-dimensional Euclidean vector space with inner product $(\cdot, \cdot)$. A \textit{root system} $\mathcal{R} \subset E$ is a finite set satisfying:
    \begin{enumerate}
        \item $\text{span}_{\mathbb{R}}(\mathcal{R}) = E$ and $0 \notin \mathcal{R}$.
        \item If $\alpha \in \mathcal{R}$ and $\lambda \alpha \in \mathcal{R}$ for $\lambda \in \mathbb{R}$, then $\lambda = \pm 1$.
        \item \textbf{Invariance under reflection:} For all $\alpha, \beta \in \mathcal{R}$, the reflection of $\beta$ along $\alpha$ is in $\mathcal{R}$:
              $$s_\alpha(\beta) := \beta - \frac{2(\alpha, \beta)}{(\alpha, \alpha)} \alpha \in \mathcal{R}$$
        \item \textbf{Crystallographic Condition:} For all $\alpha, \beta \in \mathcal{R}$, the projection is an integer:
              $$n_{\alpha\beta} := \langle \alpha, \beta \rangle = \frac{2(\alpha, \beta)}{(\alpha, \alpha)} \in \mathbb{Z}$$
    \end{enumerate}
\end{frame}

\begin{frame}[label=backup_simpleroots]{Formal Definition: Simple Roots}
    \textbf{Definition 4.4.3:}
    Let $t \in E$ be a vector such that $(t, \alpha) \neq 0$ for all $\alpha \in \mathcal{R}$.

    \begin{itemize}
        \item A root $\alpha$ is defined as \textbf{positive} ($\alpha \in \mathcal{R}^+$) if $(t, \alpha) > 0$.
        \item A positive root $\alpha \in \mathcal{R}^+$ is called \textbf{simple} if it cannot be written as the sum of two positive roots.
    \end{itemize}

    The set of simple roots is denoted by $\Delta \subset \mathcal{R}$.

    \vspace{0.3cm}
    \textbf{Theorem 3.5.5:} Every root $\alpha \in \mathcal{R}$ can be uniquely written as an integer linear combination of simple roots:
    $$ \alpha = \sum_{i=1}^r n_i \alpha^{(i)} $$
    where all $n_i \ge 0$ (if $\alpha \in \mathcal{R}^+$) or all $n_i \le 0$ (if $\alpha \in \mathcal{R}^-$).
\end{frame}

\begin{frame}[label=backup_dynkin]{Formal Definition: Dynkin Diagrams}
    \textbf{Definition 4.4.5:}
    A Dynkin diagram is a partially directed graph associated with a root system:
    \begin{itemize}
        \item The nodes are the simple roots $\alpha \in \Delta$.
        \item For $\alpha, \beta \in \Delta$ with $\|\alpha\| \le \|\beta\|$, they are connected by $\ell_{\alpha\beta}$ edges, where:
              $$ \ell_{\alpha\beta} = -n_{\alpha\beta} = 4 \cos^2 \theta_{\alpha\beta} \in \{0, 1, 2, 3\} $$
        \item If $\ell_{\alpha\beta} > 1$ (i.e. roots have different lengths), an arrow is drawn from the longer root $\beta$ to the shorter root $\alpha$.
    \end{itemize}
\end{frame}